\chapter{Hot Big Bang model}
Until now, we have studied what happens at macroscopic scales by observing the evolution of the scale factors. Thus, we now go in the opposite direction by studying the microscopic scales at which fundamental particles live and interact, and how they behave in an evolving Universe and how they contribute to its expansion. To be more precise, we are particularly interested in the very early stages of the Universe's life. 

\section{Basic concepts of Quantum Field Theory}
We now make some basic recalls of Quantum Field Theory (QFT). The main objects are the fields $\phi(x)$, which are functions of the spacetime coordinates $x^\mu$, and can be of different nature. The main different between kinds of fields is how they transform under coordinate transformations. With them, we can write a Lagrangian, more precisely a Lagrangian density, $\Lag(\phi,\p\phi)$ whose dependent on the fields, of course, and their derivatives $\p_\mu\phi(x)$. The Lagrangian needs to be invariant under general coordinate transformations, therefore it is constructed from \emph{invariants}, that have precisely this property. The integral of the Lagrangian over spacetime is the action $S$. To write a consistent action, we need to impose some conditions on the Lagrangian:
\begin{enumerate}
    \item it must be quadratic in the fields;
    \item it must have the lowest number of derivatives;
    \item it must be Lorentz-covariant.
\end{enumerate}
We now look at some examples of fields which will be important for the following discussions.

\subsubsection{Scalar field}
The simplest example of field is the \emph{scalar field} $\phi(x)$. The Lagrangian for this kind of field is 
\begin{equation}
    \Lag = X\p_\mu\phi\p^\mu\phi + Y\phi^2.
\end{equation}
Using the Euler-Lagrange equations, we can find the equations of motion for the scalar field to be 
\begin{equation}
    (X\square - Y)\phi(x) = 0.
\end{equation}
Through a Fourier transform, the field can be expressed as 
\begin{equation}
    \phi(t,\v{x}) = \int \frac{d^3k}{(2\pi)^3}e^{i(\v{k} \cdot \v{x})}\t{\phi}(t,\v{p}),
\end{equation}
where $\t{\phi}(t,\v{p})$ is the scalar field in \emph{momentum space}, or in \emph{Fourier space}. The equation of motion can be then written as 
\begin{equation}
    \p_0^2\t{\phi}(t,\v{p}) + \omega_k^2\t{\phi}(t,\v{p}) = 0,\quad \omega_k^2 \equiv c^2k^2 - \frac{Y}{X}c^2,
\end{equation}
which is just the equation of a harmonic oscillator with pulse $\omega_k$. From Quantum Mechanics, we already know that $E = \hbar\omega$ and that $\v{q} = \hbar\v{k}$, so using them with the mass-shell condition $E^2 = q^2c^2 + m^2c^4$, we can identify 
\begin{equation}
    \frac{Y}{X} = -\frac{m^2c^2}{\hbar^2}.
\end{equation}
Hence, the equation of motion written before become the so-called \emph{Klein-Gordon equation}
\begin{equation}
    \bigg(\square + \frac{m^2c^2}{\hbar^2}\bigg)\t{\phi}(t,\v{p}) = 0.
\end{equation}
\begin{obs}
    Until 2012 with the discovery of the Higgs boson, the scalar field did not have an example in nature, being just a mathematical artefact.
\end{obs}

\subsubsection{Vector field} 
The most familiar example of vector field is $A_\mu$ that appears in Electromagnetism. Maxwell equations, in covariant form, are 
\begin{equation}
    \begin{cases}
        \p_\mu F^{\mu\nu} = 0 \\
        \p_{[\lambda}F_{\mu\nu]} = 0
    \end{cases},
\end{equation}
where $F_{\mu\nu} = \p_\mu A_\nu - \p_\nu A_\mu$ is a \emph{field strength}, with $A^\mu = (\Phi/c,A^i)$ where $\Phi$ is the \emph{scalar potential}, typically $\Phi \equiv A^0$, and $A^i$ is the \emph{vector potential}. Anyway, here we can start to see some differences with respect to the scalar field case. By looking at the Maxwell equations, we can see that $A^\mu$ does not appear ``by itself'', but only in terms of its derivates. Therefore, $A^\mu$ is not uniquely defined, but under transformation $A^\mu \to A^{\prime\mu} \equiv A^\mu + \p^\mu\lambda(x)$, where $\lambda(x)$ is a scalar function, Maxwell equations are invariant, or in other words $F_{\mu\nu} \to F'_{\mu\nu} \equiv F_{\mu\nu}$. This is really important because when we try to write a Lagrangian for this field we cannot use $A^\mu A_\mu$ as we did for the scalar field since this is not invariant, which implies that also the Lagrangian would not be invariant. Thus, for what we said before, we can assume that the Lagrangian contains instead $F_{\mu\nu}F^{\mu\nu}$, which is indeed invariant; that is 
\begin{equation}
    \Lag = -\frac{1}{4}F_{\mu\nu}F^{\mu\nu}.
\end{equation}
The equations of motion that we derive from this are 
\begin{equation}
    \square A^\mu = \p^\mu(\p_\nu A^\nu).
\end{equation}
Because of the invariance of the Maxwell equations under the $A^\mu \to A^{\prime\mu} \equiv A^\mu + \p^\mu\lambda(x)$ transformation, we can impose a \emph{gauge condition} which removes one degree of freedom. There are many choices for such a condition, but we use the \emph{Lorenz gauge condition}, that is 
\begin{equation}
\label{eq: Lorenz gauge condition}
    \p_\mu A^\mu = 0.
\end{equation}
In this way, the equations of motion become simply
\begin{equation}
    \square A^\mu = 0.
\end{equation}
In momentum space, this can be written as 
\begin{equation}
    \p_0^2\t{A}^\mu(t,\v{k}) + c^2k^2\t{A}^\mu(t,\v{k}) = 0.
\end{equation}
The solution can be then be expressed as 
\begin{equation}
    \t{A}^\mu(t,\v{k}) = \epsilon^{\mu}(x)e^{-i\omega_kt},\quad \omega_k^2 = c^2k^2.
\end{equation}
Because of this, the gauge condition \eqref{eq: Lorenz gauge condition} is equivalent to asking $\epsilon_\mu k^\mu = 0$, which implies that the fields $A^\mu$ and $k^\mu$ are orthogonal to each other. However, there is more. The Lorenz gauge condition does not remove completely the gauge freedom; in fact, after a transformation $A^\mu \to A^{\prime\mu} \equiv A^\mu + \p^\mu\lambda(x)$, the Lorenz gauge condition becomes 
\begin{equation}
    \p_\mu A^{\prime\mu} = \p_\mu A^\mu + \square\lambda(x) = 0,
\end{equation}
and since $\p_\mu A^\mu = 0$, we can maintain the Lorenz gauge condition if and only if $\square \lambda(x) = 0$. This is what is called the \emph{residual gauge freedom}. This condition removes another degree of freedom, lowering them from four to two removing the non-physical degrees of freedom, while the remaining two are the physical ones, that can be interpreted as the polarization states of the photon. 

\subsubsection{Spinor field}
There are a lot of examples for spinor fields, but the most famous one is maybe the electron. Anyway, a spinor field is often indicated as $\Psi^A(x)$ where the index $A$, although $A = 0,\dots,3$, it is not a Lorentz index, instead is a spinor index. A spinor transforms under a Lorentz transformation as 
\begin{equation}
    \Psi^A \to \Psi^{\prime A} \equiv S(\Lambda)^A_B\Psi^B.
\end{equation}
We can also define the \emph{adjoint spinor} $\b{\Psi}^A$, which transforms as 
\begin{equation}
    \b{\Psi}_A \to \b{\Psi}^{\prime}_A \equiv \b{\Psi}_BS(\Lambda)^B_A.
\end{equation}
Therefore, it is clear that $\b{\Psi}_A\Psi^A$ is invariant under Lorentz transformations. Hence, the Lagrangian can be written as 
\begin{equation}
    \Lag = Xi\b{\Psi}_B\sl{\p}_A^B\Psi^B - Y\b{\Psi}_A\Psi^A,\quad \sl{\p} \equiv \gamma^\mu\p_\mu,
\end{equation}
where $\gamma^\mu$ are the Dirac matrices, which satisfy the Clifford algebra defined by the anti-commutation relation $\acomm{\gamma^\mu}{\gamma^\nu} = 2\eta^{\mu\nu}\IM$. Using the Euler-Lagrange equations with respect to the adjoint spinor, we can derive the equations of motion for the ordinary spinor 
\begin{equation}
    (Xi\sl{\p}^A_B- Y\delta^A_B)\Psi^B = 0.
\end{equation}
In the Fourier space we have 
\begin{equation}
    \{X[\omega t(\gamma^0)^A_B - k_i(\gamma^i)^A_B] - \hbar Y\delta^A_B\}\t{\Psi}^B = 0.
\end{equation}
Defining the matrix $M^A_B \equiv X[\omega t(\gamma^0)^A_B - k_i(\gamma^i)^A_B]$ and noting that $\delta^A_B \equiv \IM$,  we can write the equation of motion as 
\begin{equation}
    (M_B^A - \hbar Y\IM)\t{\Psi}^B = 0.
\end{equation}
Multiplying this last equation of the left with $(M_B^A + \hbar Y\IM)$ we find 
\begin{equation}
    (M_B^A - \hbar Y\IM)(M_B^A - \hbar Y\IM)\t{\Psi}^B = [(M_B^A)^2 - (\hbar Y)^2\IM]\Psi^B = 0.
\end{equation}
Using the Clifford algebra $\acomm{\gamma^\mu}{\gamma^\nu} = 2\eta^{\mu\nu}\IM$, we have 
\begin{equation}
    (M_B^A)^2 = X^2\gamma^\mu\gamma^\nu p_\mu p_\nu = X^2p_\mu p^\mu \IM.
\end{equation}
Thus, we remain with 
\begin{equation}
    [X^2p_\mu p^\mu - (\hbar Y)^2]\t{\Psi} = 0.
\end{equation}
For non-trivial solutions, we need to impose 
\begin{equation}
    p_\mu p^\mu = \frac{\hbar^2Y^2}{X^2},
\end{equation}
from which we find, comparing it with the relativistic relation $p_\mu p^\mu = m^2c^2$, that 
\begin{equation}
    \frac{X}{Y} = \frac{\hbar}{mc}.
\end{equation}
Hence, the equation of motion  equations of motion written before become the so-called \emph{Dirac equation}
\begin{equation}
    \bigg(\frac{i\hbar}{mc}\sl{\p}_B^A - \IM\bigg)\Psi^B = 0.
\end{equation}

Regarding the spinor field we can say more. If we perform a \emph{local transformation}, which is a transformation that is dependent on the spacetime coordinates, defined as 
\begin{equation}
    \Psi^A \to \Psi^{\prime A} \equiv \Psi^Ae^{i\theta(x)},\quad \b{\Psi}^A \to \b{\Psi}^{\prime A} \equiv \b{\Psi}^Ae^{-i\theta(x)}.
\end{equation}
The product $\b{\Psi}_A\Psi^A$ is still invariant; however, the derivative of the spinor field is not, in fact 
\begin{equation}
    \p_\mu \Psi^A \to \p_\mu\Psi^{\prime A} \equiv \p_\mu(e^{i\theta(x)}\Psi^A) = e^{i\theta(x)}[\p_\mu + i(\p_\mu\theta)]\Psi^A,
\end{equation}
which indeed breaks the Lagrangian invariance. To preserve the invariance, we need to redefine the derivative operation as 
\begin{eqnarray}
    D_\mu := \p_\mu - B_\mu,\quad B_\mu \to B_\mu' \equiv B_\mu + i\p_\mu\theta(x),
\end{eqnarray}
which is exactly what we did when we defined the covariant derivative previously. Thus, we have 
\begin{equation}
    D_\mu\Psi^A \to (D_\mu\Psi^A)e^{i\theta(x)}.
\end{equation}
The transformation with which we defined the covariant derivative has the same form of the transformation of $A_\mu$ that leaves invariant Maxwell equations. Hence, if we set $\lambda(x) \equiv (Q/\hbar)\theta(x)$, we just need to replace
\begin{equation}
    B_\mu = \frac{iQ}{\hbar}A_\mu,
\end{equation}
where $Q$ is the charge of the spinor. The covariant derivative then becomes 
\begin{equation}
    D_\mu = \p_\mu - \frac{iQ}{\hbar}A_\mu.
\end{equation}
In this way we have a locally gauge-invariant theory that describes the interaction between vector fields and spinor fields, and so, for examples, between photons and electrons. The Lagrangian is 
\begin{equation}
    \Lag = \b{\Psi}[i\gamma^\mu(\hbar\p_\mu - iQA_\mu) - mc^2]\Psi.
\end{equation}

We can also give to $D_{\mu}$ a geometric meaning. To see this, we compute the commutator between derivatives, which results in
\begin{equation}
    \comm{D_\mu}{D_\nu}\Psi^A = -\frac{iQ}{\hbar}F_{\mu\nu}\Psi^A,
\end{equation}
where the Riemann curvature tensor describes the curvature associated with the spacetime connection. By analogy, the field strength $F_{\mu\nu}$ represents the curvature associated with the gauge connection $A_{\mu}$, while $D_{\mu}$ is the corresponding gauge-covariant derivative acting on fields in the internal space.

More generally, non-Abelian gauge symmetries lead to the Yang-Mills theories underlying the other fundamental interactions. We shall not go into details: the key point is that quantized gauge fields provide the mediators of interactions. The Standard Model organizes these interactions according to the combined action of the Lie groups $\group{SU}{3}_C \times \group{SU}{2}_L \times \group{U}{1}_Y$. Its matter fields consist of quarks and leptons arranged in three generations. The $\group{SU}{3}_C$ sector describes the strong interaction and contains eight gluon fields. The electroweak sector $\group{SU}{2}_L \times \group{U}{1}_Y$ leads to the photon and to the $W^\pm$ and $Z$ bosons. Unlike the photon, the weak bosons are massive: their masses originate from the coupling to the non-zero vacuum value of the Higgs field, not from inserting an explicit gauge-breaking vector mass term in the fundamental action.

\section{Distribution functions}
To use this microscopic picture in Cosmology, we need a quantitative measure of how frequently interactions occur. The Lagrangian of an interactive theory can be factorized in the sum of the free Lagrangian $\Lag_0$ and the interactive Lagrangian $\Lag_{\rem{int}}$. The free Lagrangian specifies the propagation of the particles and their dispersion relations. The interactive part instead describes the scattering, decay, annihilation, and creation of particles. When a perturbative description of the theory is possible, its contributions can be organized using Feynman's diagrams, which are a computational representation of terms in an expansion of a \emph{transition amplitude}, not literal particle trajectories.

For a generic process $A + B \to C + D$, we denote the amplitude as $\amp^{AB}_{CD}$. The relevant probability is the given by its moduli $|\amp^{AB}_{CD}|$. Combining this with the phase-space of the final particles and the incoming flux produces the interaction \emph{cross-section} denoted as $\sigma^{AB}_{CD}$. Thus, QFT proves the microscopic quantity $\sigma(E)$ needed to describe interactions between actual populations of particles as a function of the energy $E$ of the process, For example, in processes governed by a dimensionless coupling $\alpha$ and a single characteristic relativistic energy $E$m dimensional reasoning often gives $\sigma(E) \simeq \alpha^2E^{-2}c^2\hbar^2$, however masses, long-rage forces, and afftional dynamical states can change this estimate.  

Now consider a particle $A$ moving through a population of particles $B$, with number density $n_B$. During a short time $dt$ it defines an effective interaction volume $\sigma^{AB}_{CD}v_{AB}dt$, where $v_{AB}$ is the relative speed between particles $A$ and $B$. The mean number of targets encountered is $n_B\sigma^{AB}_{CD}v_{AB}dt$, suggesting that the characteristic interaction rate is given by 
\begin{equation}
    \Gamma_{AB} = n_B\vmedio{\sigma^{AB}_{CD}v_{AB}dt}.
\end{equation}

Interaction rates tell us how many interactions happen per unit of time, for a certain kind of particle and interaction. In Cosmology, we have another timescale that we need to consider, that is the expansion rate $H$. This gives us a cosmological criterion: when $\Gamma \gg H$, many interactions can take place in a Hubble time; while when $\Gamma \ll H$, that particular process is inefficient on a cosmological timescale. This comparison alone will not guarantee that a bound state will form, but it tells us whether the microscopic interactions can keep the pace with the cosmological expansion. We have already seen that the Universe's expansion affects the particles' momenta since $|\v{q}| \simeq a^{-1}$, and consequently their energy $E(q)$. To determine $\Gamma$ as a function of cosmic time, we must know how many particles are present and how their momenta are distributed. 

It is more useful here to consider a dimensionless distribution function, factoring out the phase-space units. In particular, when we expanded the fields as plain waves we had wavenumbers $\v{k} = 2\pi\v{n}/L$ where $L$ was the physical size of the volume $V = L^3$, with $n_i \in \bb{Z}$ for each spatial direction $i$. Since for quantum particles, $\v{q} =\hbar \v{k}$, the momentum-space spacing is $2\pi\hbar$, which leads to rescaling $d^3q \mapsto d^3q/(2\pi\hbar)^3$. Moreover, quantum particles also come in different \emph{multiplicity}, in other words they have different degrees of freedom, for example spin states. Denoting the multiplicity for particles of type $s$ as $g_s$, we can write the number of quanta of the given type in an infinitesimal mass-shell volume element as 
\begin{equation}
    dN_s = \frac{g_s}{(2\pi\hbar)^3}f_s(\v{x}_p,\v{q};t)\,d^3x_p\,d^3q.
\end{equation}
In these units, the distribution function now defines the mean occupation number per one-particle state. The number density of particles of type $s$ is then given by 
\begin{equation}
\label{eq: number density of particles of type s}
    n_s(\v{x}_p) = \frac{g_s}{(2\pi\hbar)^3}\int d^3q\,f_s(\v{x}_p,\v{q}),
\end{equation} 
and therefore, integrating over the volume, we have the total number of particles 
\begin{equation}
\label{eq: number of particles of type s}
    N_s = \int_V d^3x_p\,n_s(\v{x}_p).
\end{equation}
The total energy will be simply given by the sum of the energy of each particle $E_s(q) = \sqrt{m_s^2c^4 + q^2c^2}$, where $m_s^2$ is the mass of the particle of the $s$th species and $q \equiv |\v{q}|$. Hence, the energy density is 
\begin{equation}
\label{eq: energy density of particles of type s}
    \epsilon_s(\v{x}_p) = \frac{g_s}{(2\pi\hbar)^3}\int d^3q\,f_s(\v{x}_p,q)\,E_s(q),
\end{equation}
and the mean \emph{internal energy} inside the volume is 
\begin{equation}
\label{eq: internal energy density of particles of type s}
    U_s = \int d^3x_p\,\epsilon_s(\v{x}_p).
\end{equation}
In the end, we can also find the pressure to be 
\begin{equation}
\label{eq: pressure of particles of type s}
    p_s(\v{x}_p) = \frac{g_s}{(2\pi\hbar)^3}\int d^3q\,\frac{q^2c^2}{3E_s(q)}f_s(\v{x}_p,q).
\end{equation}
All of these relations however do not determine the distribution function $f_s(\v{x}_p,q)$. The way to retrieve it is to realize that we are describing a large number of particles, but we are only interested in their collective, macroscopic behaviour. 

Consider first an isolated system which contains a single species of particles. The total energy and particle number are therefore conserved, but the particles can continuously rearrange themselves among the available microscopic states. The essential observation of statistical mechanics is that not all macroscopic distributions correspond to the same number of microscopic arrangements. The most probable distribution is the one that can be realized in the greatest number of ways. This leads us to the concept of entropy. Entropy can be defined in a number of ways, but the most fundamental one might be the Boltzmann's statistical entropy 
\begin{equation}
    S = k_B\ln{\Xi},
\end{equation}
where $\Xi$ is the number of microscopic configurations compatible with the macroscopic states. The distribution function on his side counts the number of one-particle states in the phase-space. Thus, if we coarse-grain the future mass-shell volume in cells small enough that $f$ is approximately constant within each, the number of available one-particle states in the $i$th cell is 
\begin{equation}
    G_i = \frac{g}{(2\pi\hbar)^3}\Delta^3x_{p,i}\Delta^3q_{i}.
\end{equation} 
Then, if $N_i$ particles occupy these cells, $N_i = fG_i$. The problem now is to find how many distinguishable microscopic ways we can arrange $N_i$ particles among $G_i$ states.

\subsection{Fermionic and bosonic distributions}
We start discussing fermions since they have the constraint given by Pauli's exclusion principle which tells us that a given one-particle state can contain either one or no fermion. Therefore, distributing $N_i$ in $G_i$ states is equivalent to counting how many $N_i$ states are occupied out of the $G_i$ available. The number of possibilities is then just the binomial coefficient
\begin{equation}
    \Xi_{\rem{FD}} = 
    \begin{pmatrix}
        G_i \\
        N_i    
    \end{pmatrix} = \frac{G_i!}{N_i!(G_i - N_i)!}.
\end{equation}
Thus, the entropy is 
\begin{equation}
    \begin{split}
        S_{\rem{FD}} &= k_B[\ln{(G_i!)} - \ln{(N_i!)} - \ln{(G_i - N_i)!}] \\
        &\simeq -k_BG_i[f\ln{(f)} - (1 - f)\ln{(1 - f)}],
    \end{split}
\end{equation}
where we have used Stirling's approximation for the logarithm of a factorial $\ln{n!} \simeq n\ln{(n)} - n$ for $n \gg 1$. Also, in the continuous limit $\Delta^3 \mapsto d^3$, hence 
\begin{equation}
    S[f_{\rem{FD}}] = -\frac{g_sK_B}{(2\pi\hbar)^3}\int d^3x_p\,d^3q\,[f_{\rem{FD}}\ln{(f_{\rem{FD}})} + (1 - f_{\rem{FD}})\ln{(1 - f_{\rem{FD}})}],
\end{equation}
where we emphasize this is not a general formula, but it holds for fermions which obey the Fermi-Dirac distribution, hence the ``FD'' subscript.

Thermodynamic equilibrium is the realized by the distribution function that maximizes entropy. At equilibrium, the first law of Thermodynamics relates infinitesimal variations of internal energy, entropy, volume, and particle number through 
\begin{equation}
\label{eq: first law of Thermodynamics}
    dU = T\,dS - p\,dV + \mu\,dN,
\end{equation}
where $T$ is the temperature and $\mu$ is the chemical potential.  Since we are seeking precisely the equilibrium distribution, we can impose this relation on the corresponding statistical functionals. Using \eqref{eq: first law of Thermodynamics}, we can impose the stationarity condition $\delta(U - TS + pV - \mu N_i)|_{\rem{eq}} = 0$, which implies 
\begin{equation}
    \delta S = \frac{1}{T}\delta U + \frac{p}{T}\delta V - \frac{\mu}{T}\delta N_i,
\end{equation}
hence 
\begin{multline}
    -\frac{g k_B}{(2\pi\hbar)^3}\int d^3x_p\,d^3q\,[\ln{(f_{\rem{FD}})} - \ln{(1 - f_{\rem{FD}})}]\delta f_{\rem{FD}} = \\
     = \frac{g}{T(2\pi\hbar)^3}\int d^3x_p\,d^3q\,(E - \mu)\delta f_{\rem{FD}},
\end{multline}
where we have noticed that the volume does not depend on the distribution function, whereas the number of particles and internal energy within $V$ do. 

Since the variations are arbitrary, by solving for $f_{\rem{FD}}$ we eventually obtain the \emph{Fermi-Dirac distribution}
\begin{equation}
\label{eq: Fermi-Dirac distribution}
    f_{\rem{FD}} = \frac{1}{e^{(E - \mu)/k_BT} + 1}.
\end{equation}
An analogous reasoning, but for the combinatorics, can be made for bosons, leading us to the \emph{Bose-Einstein distribution}
\begin{equation}
\label{eq: Bose-Einstein distribution}
    f_{\rem{BE}} = \frac{1}{e^{(E - \mu)/k_BT} - 1}.
\end{equation}
Then, when the mean occupation number is very small, namely when $f = N_i/G_i \ll 1$, multiple occupation of an individual quantum state is negligible and the counting simplifies to $\Xi_{\rem{MB}} \simeq G_i^N/N_i!$. This corresponds to the \emph{Maxwell-Boltzmann distribution}
\begin{equation}
\label{eq: Maxwell-Boltzmann distribution}
    f_{\rem{MB}} = e^{-(E - \mu)/k_BT}.
\end{equation}
Note that this is valid in general. As long as we work in the reference frame of the system, the
particle energy depend only on momentum via the mass-shell, whereas different regions
can be in local thermodynamic equilibrium, with different temperatures and chemical potentials. In a homogeneous Universe however, $f_s(\v{x}_p,\v{q}) \equiv \b{f}_s(q)$. This implies that we can write $\b{N}_s = V\b{n}_s$ and $\b{U}_s = V\b{\epsilon}_s$. Furthermore, the momentum-space integrals can be simplified with $d^3q \mapsto (4\pi q^2)dq$. Hence, \eqref{eq: number density of particles of type s}, \eqref{eq: energy density of particles of type s}, and \eqref{eq: pressure of particles of type s} can be written as 
\begin{equation}
\label{eq: number density of particles of type s (HU)}
    \b{n}_s(T) = \frac{g_s}{2\pi^2\hbar^3}\int_0^\infty dq\,\frac{q^2}{e^{[E_s(q) - \mu_s]/k_BT} \pm 1},
\end{equation}
\begin{equation}
\label{eq: energy density of particles of type s (HU)}
    \b{\epsilon}_s(T) = \frac{g_s}{2\pi^2\hbar^3}\int_0^\infty dq\,\frac{q^2E_s(q)}{e^{[E_s(q) - \mu_s]/k_BT} \pm 1},
\end{equation}
\begin{equation}
\label{eq: pressure of particles of type s (HU)}
    \b{p}_s(T) = \frac{g_s}{2\pi^2\hbar^3}\int_0^\infty dq\,\frac{q^4c^2}{3E_s(q)(e^{[E_s(q) - \mu_s]/k_BT} \pm 1)}.
\end{equation}
Performing the change of variables $X_s = m_sc^2/k_BT$, $W_s = \mu_s/k_BT$, and $Y = qc/k_BT$, we can write the first two relations like 
\begin{equation}
    \b{n}_s(T) = \frac{k_B^3}{2\pi^2\hbar^3c^3}g_sT^3I_\pm(X_s,W_s),
\end{equation}
\begin{equation}
    \b{\epsilon}_s(T) = \frac{k_B^4}{2\pi^2\hbar^3c^3}g_sT^4J_\pm(X_s,W_s),
\end{equation}
where we have also defined the integral functions 
\begin{equation}
\label{eq: I+- integral function}
    I_\pm(X_s,W_s) \equiv \int_0^\infty dY\,\frac{Y^2}{\exp{\{\sqrt{X_s^2 + Y^2} - W_s\}} \pm 1},
\end{equation}
\begin{equation}
\label{eq: J+- integral function}
    J_\pm(X_s,W_s) \equiv \int_0^\infty dY\,\frac{Y^2\sqrt{X_s^2 + Y^2}}{\exp{\{\sqrt{X_s^2 + Y^2} - W_s\}} \pm 1},
\end{equation}
where the subscript ``$+$'' refers to fermions, while ``$-$'' refers for bosons.

\subsubsection{Non-relativistic particles}
In the non-relativistic limit, the momentum is smaller than the rest mass. This means that the factor $\pm1$ that appears in the distribution functions \eqref{eq: Fermi-Dirac distribution} and \eqref{eq: Bose-Einstein distribution} is negligible compared to the exponential, therefore both of them reduce to the Maxwell-Boltzmann distribution function \eqref{eq: Maxwell-Boltzmann distribution}. Moreover, if the momentum is small, we can expand in Taylor series the argument of the exponential around $Y^2 = 0$, at the second order of expansion we have $\sqrt{X^2_s + Y^2} = X_s + Y^2/(2X_s) + o(Y^4)$. Hence, the integral functions become 
\begin{equation}
    I_\pm(X_s,W_s) \simeq  e^{W_s - X_s}\int_0^\infty dY\,e^{-Y^2/(2X_s)}Y^2,
\end{equation}
\begin{equation}
    J_\pm(X_s,W_s) \simeq  e^{W_s - X_s}\int_0^\infty dY\,e^{-Y^2/(2X_s)}Y^2\bigg(x_s + \frac{Y^2}{2X_s}\bigg).
\end{equation}
Setting $A equiv 1/(2X_s)$, we remain with Gaussian integrals that can be computed, they are 
\begin{equation}
    \int_0^\infty dY\,Y^2e^{-AY^2} = -\frac{d}{dA}\int_0^\infty dY\,e^{-AY^2} = -\frac{d}{dA}\bigg(\frac{1}{2}\sqrt{\frac{\pi}{A}}) = \frac{\sqrt{\pi}}{4}A^{-3/2},
\end{equation}
\begin{equation}
    \int_0^\infty dY\,Y^4e^{-AY^2} = -\frac{d^2}{dA^2}\int_0^\infty dY\,e^{-AY^2} = -\frac{d^2}{dA^2}\bigg(\frac{1}{2}\sqrt{\frac{\pi}{A}}) = \frac{3\sqrt{\pi}}{8}A^{-5/2}.
\end{equation}
Hence, we have 
\begin{equation}
\label{eq: number density of particles of type s (HU + NRL)}
    \b{n}_s(T) = g_s\lambda_s^{-3}\exp{\bigg\{-\frac{m_sc^2 - \mu_s}{k_BT}\bigg\}},
\end{equation}
\begin{equation}
\label{eq: energy density of particles of type s (HU + NRL)}
    \b{\epsilon}_s(T) = n_s(T)m_sc^2 + \frac{3}{2}n_s(T)k_BT.
\end{equation}
We have also defined the quantity $\lambda_s$, which is called the \emph{thermal de Broglie wavelength} for particles of species $s$, as 
\begin{equation}
    \lambda_s \equiv \hbar\sqrt{\frac{2\pi}{m_sk_BT}}.
\end{equation}
These relations give us valuable insight on the behaviour of the particle gas of species $s$ in the non-relativistic limit. In equation \eqref{eq: number density of particles of type s (HU + NRL)}, we see that, if $k_BT > m_sc^2$, non-relativistic particles scale as a power law in temperature $T^{3/2}$, meaning that the higher the temperature, the denser number of particles in the volume. At lower temperatures, if the number of particles is somewhat conserved, it is usually the chemical potential that balances out the energy in the exponential, keeping the number density approximately constant. Oppositely, if no process preserves particle number, the number density becomes exponentially suppressed.

Concerning the energy density, the Taylor expansion for small momenta clearly shows us
two regimes: the lowest-energy one, where $E_s \simeq m_sc^2$; and a small correction, proportional to $q^2/(2 m_s)$. In equation \eqref{eq: energy density of particles of type s (HU + NRL)}, this corresponds, again, to having: a regime for $k_BT \ll m_sc^2$, where the energy density is just rest-mass energy weighted by the number density of particles; and another regime, for $k_BT > m_sc^2$, where the energy density is mainly due to the particle mean kinetic energy.

Finally, we recall the expression of the pressure in \eqref{eq: pressure of particles of type s}. The Taylor expansion for the small momenta give $Y^2/\sqrt{X_s^2 + Y^2} = Y^2/X_s + o(Y^3)$, thus 
\begin{equation}
\label{eq: pressure of particles of type s (HU + NRL)}
    \b{p}_s(T) = 0 + \b{n}_s(T)k_BT.
\end{equation}
In other words, the lowest-energy limit corresponds to zero pressure, like we already derived for matter. Then, when $k_BT > m_sc^2$, we find the equation of the perfect gas $pV = k_BNT$.

\subsubsection{Ultra-relativistic particles}
In the ultra-relativistic limit, the particles rest mass is either zero, or it is negligible compared to the momentum, namely $E(q) \simeq qc$. On the other hand, chemical potentials are determined by the conserved charges and the reactions maintaining chemical equilibrium. Reactions that changes the number of photons enforce $\mu_\gamma = 0$, while the tiny asymmetry between particles and antiparticles (assuming a similarly small lepton asymmetry) allows us to take $|\mu_s| \ll k_BT$ for electrons, positrons, and neutrinos. Therefore, in this regime $X_s \simeq W_s \ll Y$, and we can derive the scaling relations
\begin{equation}
    \b{n}_s(T) = \frac{k_B^3}{2\pi^2\hbar^3c^3}g_sT^3I_\pm(0,0),
\end{equation}
\begin{equation}
    \b{\epsilon}_s(T) = \frac{k_B^4}{2\pi^2\hbar^3c^3}g_sT^4J_\pm(0,0),
\end{equation}
\begin{equation}
    \b{p}_s(T) = \frac{1}{3}\b{\epsilon}_s(T).
\end{equation}
The integral functions $I_\pm(0,0)$ and $J_\pm(0,0)$ are in fact special cases of the integral 
\begin{equation}
    \int_0^\infty dx\,\frac{x^\alpha}{e^x + \beta} = \bigg(1 + \frac{1 + \beta}{2^{\alpha + 1}}\bigg)\alpha!\zeta(\alpha + 1),\quad \beta = \pm 1,\,\alpha \in \bb{N}_0,
\end{equation}
and where $\zeta(x)$ is the \emph{Riemann zeta function}. Hence, we have 
\begin{equation}
    \b{n}_s(T) = \frac{\zeta(3)k_B^3}{\pi^2\hbar^3c^3}g_sT^3\times 
    \begin{cases}
        1, &s \equiv \rem{BE} \\
        3/2, &s \equiv \rem{FD}
    \end{cases},
\end{equation}
\begin{equation}
    \b{\epsilon}_s(T) = \frac{\pi^2k_B^4}{30\hbar^3c^3}g_sT^4\times 
    \begin{cases}
        1, &s \equiv \rem{BE} \\
        7/8, &s \equiv \rem{FD}
    \end{cases},
\end{equation}
where we have used that $\zeta(4) = \pi^4/90$, while $\zeta(3) \simeq 1.2021$. 

\section{Multi-component system (WIP)}
For a joint system of particle species that are in equilibrium by themselves, the total internal energy and the total entropy are given by simply $U = \sum_sU_s$ and $S = \sum_sS_s$. Using the first law of Thermodynamics, we have that the total entropy is related to thermodynamical quantities through
\begin{equation}
    dS = \nsum_s\bigg(\frac{dU_s}{T_s} + \frac{\b{p}_s}{T_s}\,dV - \frac{\mu_s}{T_s}\,dN_s\bigg).
\end{equation}
We can consider the $dV$ term negligible since the rate of expansion and contraction of the volume is irrelevant compared to the rate of interactions that occur in the system. The only way to  have exchange of energy and momenta between species while keeping entropy at a maximum is that the temperature $T_s$ is the same for all the species.

Consider the scattering process $A + B \to A + B$. To keep entropy at a maximum, it is required that $dU_A = -dU_B$, this implies that 
\begin{equation}
    dS = \bigg(\frac{1}{T_A} - \frac{1}{T_B}\bigg)dU_A = 0,
\end{equation}
if and only if $T_A = T_B$ since $dU$ can be positive or negative. Consider now a scattering process that changes the particle's species, so for example $A + B \to C + D$. To have thermodynamical equilibrium, now it is also required that $\mu_A + \mu_B = \mu_C + \mu_D$, so that they are in chemical equilibrium, in addition to being in thermodynamical equilibrium. 

Recall that the energy density of every species evolves in time with \eqref{eq: cosmological continuity equation in GR}, and this also holds for the total energy density. Moreover, multiplying by the volume $V$ equation \eqref{eq: cosmological continuity equation in GR} and noting that 
\begin{equation}
    \frac{d\ln{V}}{dt} = 3H,
\end{equation}
we can see that $3HV = \dot{V}$ and write 
\begin{equation}
    \frac{d\b{U}}{dt} + \b{p}\frac{dV}{dt} = 0,\quad d\b{U} \equiv \b{\epsilon}V,
\end{equation}
which implies that 
\begin{equation}
    d\b{U} + \b{p}\,dV = 0.
\end{equation}
Therefore, the total entropy is 
\begin{equation}
    dS = \nsum_s\frac{\mu_s}{T_s}\,dN_s.
\end{equation}
This means that for a collection of particles that are in thermodynamical equilibrium the expansion is adiabatic. Moreover, the entropy is constant, thus we can write it as the entropy density $s$ times the volume $V$, namely $S = sV$, thus $s \propto a^{-3}$. Then, considering again the first law of Thermodynamics, written as 
\begin{equation}
    d\b{U} = T\,dS - \b{p}\,dV + \mu\,d\b{N},
\end{equation}
takes us to 
\begin{equation}
    \b{\epsilon} = Ts - \b{p} + \nsum_s \b{\mu}_s\b{n}_s.
\end{equation}
... 